% généré par scripts/06_textes.py
\dosmeta{Année de classement : 397 (directe : nom consulaire lu). Références : ICUR I 2812 = ILCV 1508 = ILMN-01, 00447. Rattachement selon l’EDB : \emph{Incertae originis} (le volume I des ICVR réunit les inscriptions dont la provenance est perdue). Conservation selon l’EDB : \emph{Napoli   ; Museo Archeologico Nazionale}. Support : Plaque de marbre. 23,5 × 36 × 4 cm (h. × l. × ép.). Inscription gravée. Reproductions signalées par l’EDB : reproduction, Camodeca et alii 2000 (droits réservés ; non reproduites). Fiche EDB : Filippo Piazzolla, première version 2015. Notice du catalogue p.~\pageref{not:405}.}
\txtbloc{Texte transmis (EDCS, état 2022)}{\lnum{1}Bonifae bene qu(i)\{a\}es\par
\lnum{2}⸢c⸣enti catec(h)umi\par
\lnum{3}no qui vixit annos\par
\lnum{4}p(lus) m(inus) LX d(epositus) IIII Kal(endas)\par
\lnum{5}Iulias Caesario et Attico\par
\lnum{6}conss(ulibus)}
\txtbloc{Texte de l’EDB (EDB32736)}{\lnum{1}Bonifae bene quaes/quenti catecumi/no qui vixit annos\par
\lnum{2}p(lus) m(inus) LX[- - -] 《d(epositus)》 〚- - -IV- - -〛 IIII ka⌜l⌝(endas)\par
\lnum{3}iulias Caesario et  Attico\par
\lnum{4}cons(ulibu)s}
